\apendice{Plan de Proyecto Software}

\section{Introducción}

Este apartado describe la planificación y organización seguida durante el desarrollo del proyecto.

El sistema desarrollado está formado por dos addons complementarios para Blender:

\begin{itemize}
    \item Un addon de recopilación de datos, encargado de registrar automáticamente la actividad del usuario durante el modelado 3D.
    
   \item Un addon de análisis y visualización, encargado de estudiar los datos recogidos y mostrar gráficos y representaciones visuales para entender mejor cómo trabaja el usuario dentro de Blender.
\end{itemize}

El proyecto se desarrolló poco a poco, creando primero versiones simples que funcionaban correctamente y mejorándolas después con nuevas funciones y correcciones.

De esa forma, fue posible identificar los errores desde el principio y poder mejorarlo con un método más organizado y seguro.

Para manejarlo todo, utilizamos GitHub, una plataforma que permite hacer rastros de las diferentes versiones del proyecto, controlar los cambios realizados e incorporar la evolución del trabajo en la plataforma a través del tiempo.

\section{Planificación temporal}

El desarrollo se organizó en sprints quincenales basados en la evolución real del proyecto.

Cada sprint incorporó nuevas funcionalidades y mejoras sobre las versiones anteriores, permitiendo evolucionar progresivamente desde un sistema básico de captura de datos hasta un entorno completo de análisis y visualización dentro de Blender.

\begin{itemize}

    \item \textbf{Fase previa (Noviembre)}  
    Estudio inicial de viabilidad y análisis de Blender. Durante esta fase se realizaron pruebas sobre la API \texttt{bpy} y se evaluó la posibilidad de capturar eventos relacionados con el modelado 3D.

    \item \textbf{Sprint 0 (Enero)}  
    Definición de requisitos, análisis del problema y diseño inicial del sistema.

    \item \textbf{Sprint 1 (Febrero - semanas 1-2)}  
    Desarrollo de la primera versión funcional del addon de captura. Implementación inicial de handlers y exportación de datos a CSV.  
    \textit{Resultado: versión inicial funcional del sistema de recopilación (v0.1).}

    \item \textbf{Sprint 2 (Febrero - semanas 3-4)}  
    Mejora de la captura de eventos y reorganización de la estructura de datos almacenada en CSV.

    \item \textbf{Sprint 3 (Marzo - semanas 1-2)}  
    Desarrollo inicial del addon de análisis y primeras pruebas de representación visual.  
    \textit{Resultado: integración inicial del sistema de análisis (v0.2).}

    \item \textbf{Sprint 4 (Marzo - semanas 3-4)}  
    Implementación de métricas temporales, espaciales y estructurales. Desarrollo de las primeras visualizaciones gráficas.

    \item \textbf{Sprint 5 (Abril - semanas 1-2)}  
    Refactorización importante del sistema de captura y mejora del sistema de logging.  
    \textit{Resultado: versión avanzada del sistema de recopilación (v0.4).}

    Durante este sprint se incorporaron varias mejoras relevantes:

    \begin{itemize}
        \item soporte completo para \textit{Edit Mode} mediante \textit{bmesh},
        \item mejora en la detección de operadores como Shift+D, Alt+D, Ctrl+V y Merge,
        \item detección robusta de coordenadas UV,
        \item optimización del sistema de logging mediante registro basado en cambios,
        \item incorporación de controles Start/Stop,
        \item implementación de paneles de depuración,
        \item y reorganización del sistema de generación de archivos CSV.
    \end{itemize}

    \item \textbf{Sprint 6 (Abril - semanas 3-4)}  
    Desarrollo avanzado del addon de análisis. Integración de gráficos 2D, visualizaciones 3D y generación procedural de geometría dentro de Blender.

    \item \textbf{Sprint 7 (Mayo - semanas 1-2)}  
    Pruebas funcionales, optimización general del sistema y corrección de errores detectados durante las validaciones.

    \item \textbf{Sprint 8 (Mayo - semanas 3-4)}  
    Redacción de documentación técnica, revisión del proyecto y preparación de la entrega final.  
    \textit{Resultado: versión final estable del sistema completo (v1.0).}

\end{itemize}

La planificación descrita se corresponde con la evolución real registrada en el repositorio GitHub del proyecto.

\subsection{Resumen de sprints}

La Tabla~\ref{tab:sprints_resumen} resume las principales fases de desarrollo realizadas durante el proyecto.

\begin{table}[H]
\centering
\small
\begin{tabular}{|c|c|p{5cm}|p{4cm}|}
\hline
\textbf{Sprint} & \textbf{Periodo} & \textbf{Funcionalidades principales} & \textbf{Resultado} \\
\hline

Sprint 1 & Febrero (1-2) & Captura inicial de eventos y exportación CSV & Versión inicial (v0.1) \\
\hline

Sprint 2 & Febrero (3-4) & Mejora de captura y organización de datos & Sistema más estable \\
\hline

Sprint 3 & Marzo (1-2) & Desarrollo inicial del addon de análisis & Integración análisis (v0.2) \\
\hline

Sprint 4 & Marzo (3-4) & Métricas y primeras visualizaciones & Sistema analítico inicial \\
\hline

Sprint 5 & Abril (1-2) & Refactorización y optimización del logger & Versión avanzada (v0.4) \\
\hline

Sprint 6 & Abril (3-4) & Visualizaciones 2D/3D y geometría procedural & Sistema integrado \\
\hline

Sprint 7 & Mayo (1-2) & Validación y corrección de errores & Sistema validado \\
\hline

Sprint 8 & Mayo (3-4) & Documentación y entrega final & Versión final (v1.0) \\
\hline

\end{tabular}
\caption{Resumen de sprints del proyecto.}
\label{tab:sprints_resumen}
\end{table}

\section{Viabilidad del proyecto}

\subsection{Viabilidad técnica}

El proyecto es técnicamente viable gracias a las capacidades que ofrece Blender mediante su API en Python.

El sistema desarrollado utiliza handlers, operadores personalizados y acceso directo a estructuras internas de Blender para registrar información relacionada con la interacción del usuario y el estado de la escena.

Además, el segundo addon permite procesar posteriormente dichos datos mediante:

\begin{itemize}
    \item cálculo de métricas,
    \item generación de gráficos 2D,
    \item visualizaciones tridimensionales,
    \item y representación procedural de información dentro de Blender.
\end{itemize}

La separación entre captura y análisis permitió desarrollar un sistema modular y desacoplado, facilitando tanto el mantenimiento como la ampliación futura del proyecto.

\subsection{Viabilidad económica}

El proyecto se desarrolló utilizando principalmente herramientas de software libre, por lo que los costes económicos directos fueron reducidos.

\subsubsection{Costes de software}

\begin{itemize}
    \item Blender: 0 €
    \item Python: 0 €
    \item GitHub: 0 €
    \item Librerías utilizadas: 0 €
\end{itemize}

\subsubsection{Costes de hardware}

Se utilizó un equipo personal con un valor aproximado de 1200 €. Considerando una amortización estimada de cuatro años y una duración del proyecto de cuatro meses:

\[
\text{Coste hardware} = \frac{1200}{48} \times 4 = 100 \,€
\]

\subsubsection{Coste estimado de desarrollo}

Tomando como referencia un perfil junior de desarrollo software:

\[
1200 \times 4 = 4800 \,€
\]

\subsubsection{Coste total estimado}

\[
100 + 4800 = 4900 \,€
\]

Aunque el proyecto se desarrolló en un entorno académico, su equivalente aproximado en un contexto profesional sería cercano a los 4900 €.

\subsection{Viabilidad legal}

Blender se distribuye bajo licencia GPL, por lo que los addons desarrollados mantienen compatibilidad con dicha licencia.

Asimismo, las librerías utilizadas son compatibles con un uso académico y redistribuible.

El sistema incorpora mecanismos de consentimiento para informar al usuario sobre la recopilación de datos antes de iniciar la monitorización.

\section{Análisis de riesgos}

Uno de los principales riesgos del proyecto fue la complejidad de la API de Blender, especialmente en la detección de eventos dentro de distintos modos de trabajo como \textit{Object Mode} y \textit{Edit Mode}.

También aparecieron dificultades relacionadas con:

\begin{itemize}
    \item la detección de operadores complejos,
    \item la captura de coordenadas UV,
    \item la generación de grandes volúmenes de datos,
    \item y la integración entre ambos addons.
\end{itemize}

Para reducir estos riesgos se aplicaron distintas estrategias:

\begin{itemize}
    \item desarrollo iterativo,
    \item validación continua,
    \item refactorización progresiva,
    \item optimización del sistema de logging,
    \item y uso de una arquitectura modular desacoplada.
\end{itemize}

Estas medidas permitieron mejorar progresivamente la estabilidad y fiabilidad del sistema durante el desarrollo.